\begin{tikzpicture}
\node[minimum width=15cm,minimum height=9cm] (current page) {};
\begin{scope}[overlay]

\tikzset{
    box/.style={
        align=left,
        draw=red,
        fill=white,
        ultra thick,
        anchor=north,
        at={([yshift=-1cm]current page.center)},
    },
}
    \begin{visibleenv}<2>
    \node[box] {
        p4info\_helper, switch objects created by setup code
    };
    \end{visibleenv}
    \begin{visibleenv}<3>
    \node[box] {
        full name of table, including stage it is defined in
    };
    \end{visibleenv}
    \begin{visibleenv}<4>
    \node[box] {
        match value --- format would be different \\
        if key was \texttt{lpm} or \texttt{ternary} \\
        instead of exact match
    };
    \end{visibleenv}
    \begin{visibleenv}<5>
    \node[box] {
        \texttt{write\_or\_overwrite\_table\_entry} not the `raw' function \\
        (one I wrote based on one P4 tutorial authors wrote) \\
        uses gRPC (remote procedure call) library, which adds some extra steps
    };
    \end{visibleenv}
\end{scope}
\end{tikzpicture}
